\chapter{Conclusions}
\label{chap:conclusions}

Local robotic polishing requires a strategy for movement, contact loading, and time within a selected region. Human region specification supplies the spatial task, but does not by itself determine that temporal strategy. This thesis has addressed the resulting skill-generation problem through paired motion--force teaching, observation-conditioned sequence learning, and a common robot execution interface.

Tool position is learned relative to a frozen stain center and paired with orientation and target force in one sequence. The Flow Matching policy uses RGB, Cartesian state, and three-axis force history. Its direction adapter aligns live XY and the policy image, then restores generated XY by inverse rotation and center addition; Z, orientation, and force keep their conventions. E1/E2 disable direction alignment; E3 reuses the same C checkpoint and normalizer for registered planar transfer. The interfaces separate skill generation, contact control, and completion recognition.

The first research question compares strategy providers for the fixed human-specified task. In the primary E1 cohort, rules, force-inclusive replay, and the conditioned policy each complete the required contact--loaded-pass--exit procedure in all five planned executions. Replay has a shorter mean completion time and smaller time dispersion than C, while C has greater mean excess-force exposure. The learned skill is therefore an executable strategy through the common robot interface, and rules and replay are also effective alternatives within this nominal task. C5 remains in the original C comparison and C6 is a separate diagnostic. The strategy summary distinguishes additional \(R^{*}\) from original T/C. Processing retiming is associated with a \RstarTimeGain~s (\RstarTimePercent\%) shorter mean than original R while preserving +18 N and the common execution settings. This separate diagnostic uses a quintic peak-feed parameter and does not establish an IL advantage at matched measured speed.

The second research question compares the force-information configurations while retaining lower-level feedback. The force-observing motion--force policy C has a \BCSeconds~s (\BCPercent\%) shorter mean operator-based completion interval than force-masked B, with five completed attempts in each. Motion-only A completes four of five attempts; A4 does not progress to release and is interrupted. It remains in all-attempt force, travel, and elapsed-interval statistics, while A's successful-time mean uses four completions. Timing, phase progression, and force exposure differ across the configurations, without a monotonic improvement from A through B to C or an identified cause for the single interrupted transition.

The third research question evaluates spatial reuse through an image-based processing outcome. In \EthreeMainN{} selected N/P1/D1/D2 runs, \EthreeImageN{} pairs permit whole-stain evaluation. Mean \(\pm\) sample SD board-marker area reductions are \EthreeNRemoval\(\pm\)\EthreeNSD, \EthreePoneRemoval\(\pm\)\EthreePoneSD, \EthreeDoneRemoval\(\pm\)\EthreeDoneSD, and \EthreeDtwoRemoval\(\pm\)\EthreeDtwoSD\%, respectively. D1-4 remains selected with an unavailable removal value due to tool occlusion. P2's five executed, occluded runs remain in the 25-run archive outside the main removal comparison. The results support reuse of the same policy together with task registration and the direction adapter within the tested planar placements, without new teaching or retraining.

The demonstrated scope combines contact-skill execution, observed force-information differences, and E3 image-based ink-area reduction. The latter neither establishes E1/E2 removal quality nor independently measures metal surface restoration, roughness, or micro-scratch removal. It also does not establish arbitrary geometry or height robustness. Existing dose, blending, and adaptive-execution assets remain separate extensions with their documented status.

Together, the three evaluations establish executable learned motion--force skills, observed force-information differences, and registered planar reuse with image-area evidence. Rules and replay retain their demonstrated value for the fixed task.
